% Generated by scripts/generate_exercise_bank.py; edit book/exercises.json.
\subsection{按章练习与累积复习参考答案}
\phantomsection\label{app:answers}
先遮住解题过程，只读第一步；独立完成后再核对结果、单位与形状。各组标题给出题目起始页，电子版点击“第一步”可回到原题。
\Needspace{6\baselineskip}
\subsubsection*{序章\enspace{\mdseries\footnotesize（\hyperref[exercise:C-001]{题目见第~\pageref*{exercise:C-001}~页起}）}}
\begin{enumerate}[label={},leftmargin=4.8em,itemsep=4pt]
\item[\textbf{C-001.}] \phantomsection\label{solution:C-001}\label{solution:verify-generated-news}\hyperref[exercise:verify-generated-news]{\textbf{第一步：}}先找带日期的原始出处。\par\nobreak\textbf{解题过程：}检查带时间戳的权威或原始来源，交叉核验具体数字与上下文；Prompt 不能代替外部证据。
\end{enumerate}
\Needspace{6\baselineskip}
\subsubsection*{第~\ref*{intro:00-functions}~章 · 从可调函数到损失：怎样比较预测？\enspace{\mdseries\footnotesize（\hyperref[exercise:C-002]{题目见第~\pageref*{exercise:C-002}~页起}）}}
\begin{enumerate}[label={},leftmargin=4.8em,itemsep=4pt]
\item[\textbf{C-002.}] \phantomsection\label{solution:C-002}\label{solution:linear-prediction}\hyperref[exercise:linear-prediction]{\textbf{第一步：}}先算斜率项 $wx$：$x$ 为负时 $2x$ 也为负，再加偏置 $1$。\par\nobreak
\textbf{解答：} $(-1,1,5)$。
\item[\textbf{C-003.}] \phantomsection\label{solution:C-003}\label{solution:target-change}\hyperref[exercise:target-change]{\textbf{第一步：}}先检查变化发生在输入还是目标。\par\nobreak\textbf{解题过程：}要学习变化后的目标，应收集相应训练或校准样本并更新参数／模型规则；仅更换测试输入不能解决同一 $x$ 对应新标签的问题。
\item[\textbf{C-004.}] \phantomsection\label{solution:C-004}\label{solution:three-point-mse}\hyperref[exercise:three-point-mse]{\textbf{第一步：}}先列出三项残差，再平方求平均。\par\nobreak
\textbf{解答：} 预测为 $(-1,0,1)$、真实值为 $(-1,2,2)$，故 $\mathrm{MSE}=(0^2+(-2)^2+(-1)^2)/3=5/3$。
\item[\textbf{C-005.}] \phantomsection\label{solution:C-005}\label{solution:bias-candidate-comparison}\hyperref[exercise:bias-candidate-comparison]{\textbf{第一步：}}把两组参数的三个残差分别列全。\par\nobreak\textbf{解题过程：}$b=1$ 时残差为 $(1,-1,0)$，均方误差 $2/3$；$b=2$ 时残差为 $(2,0,1)$，均方误差 $5/3$。单个样本猜对不等于整体损失更小，这两个候选中应选 $b=1$。
\item[\textbf{C-006.}] \phantomsection\label{solution:C-006}\label{solution:bias-shift}\hyperref[exercise:bias-shift]{\textbf{第一步：}}先把每项旧预测加 $1$。\par\nobreak
\textbf{解题过程：}新预测为 $\answerresult{(0,1,2)}$，残差从 $(0,-2,-1)$ 变为 $(1,-1,0)$。第一个样本更远离真实值，第二、三个更接近，第三个已猜对。
\end{enumerate}
\Needspace{6\baselineskip}
\subsubsection*{第~\ref*{intro:16a-gradient-descent}~章 · 怎样让损失变小？梯度下降调参数\enspace{\mdseries\footnotesize（\hyperref[exercise:C-007]{题目见第~\pageref*{exercise:C-007}~页起}）}}
\begin{enumerate}[label={},leftmargin=4.8em,itemsep=4pt]
\item[\textbf{C-007.}] \phantomsection\label{solution:C-007}\label{solution:local-derivative-change}\hyperref[exercise:local-derivative-change]{\textbf{第一步：}}先求导，再用“导数乘参数的微小变化”估计损失变化。\par\nobreak
\textbf{解题过程：}$L'(w)=w-3$，所以 $L'(1)=\answerresult{-2}$，估计 $\Delta L\approx(-2)\times0.1=\answerresult{-0.2}$。直接计算得 $L(1)=2$、$L(1.1)=1.805$，实际变化 $-0.195$。导数描述当前位置的变化率，而这段区间内的斜率也在变；缩小步长时，局部估计更准确。
\item[\textbf{C-008.}] \phantomsection\label{solution:C-008}\label{solution:synchronous-gradient-step}\hyperref[exercise:synchronous-gradient-step]{\textbf{第一步：}}两项梯度都在旧参数 $(1,1)$ 处计算。\par\nobreak
\textbf{解题过程：}梯度为 $(-2/3,0)$，故 $w_{\rm new}=1-0.3\times(-2/3)=1.2$、$b_{\rm new}=1$，新参数为 $\answerresult{(1.2,1)}$。损失由 $2/3\approx0.667$ 降为 $\tfrac23(1.2-1.5)^2+\tfrac12=\answerresult{0.56}$。$b$ 的梯度为零，所以这一步只改变 $w$。
\item[\textbf{C-009.}] \phantomsection\label{solution:C-009}\label{solution:mini-batch-gradient-average}\hyperref[exercise:mini-batch-gradient-average]{\textbf{第一步：}}先固定 $(w,b)=(1,0)$，算完两份梯度再取平均。\par\nobreak
\textbf{解题过程：}两项残差分别为 $0,-1$，样本梯度为 $(0,0)$、$(-2,-2)$，平均为 $\answerresult{(-1,-1)}$。同步更新得到 $(w,b)=\answerresult{(1.1,0.1)}$。计算批次内第二份梯度时仍用旧参数；逐样本更新会得到另一种训练过程。
\item[\textbf{C-010.}] \phantomsection\label{solution:C-010}\label{solution:quadratic-gradient-transfer}\hyperref[exercise:quadratic-gradient-transfer]{\textbf{第一步：}}先写出 $\partial L/\partial w=\tfrac43w-2$、$\partial L/\partial b=2b-2$，再求每步参数偏差的乘数。\par\nobreak
\textbf{解题过程：}$\nabla L(0,3)=(-2,4)$；以 $\eta=0.3$ 同步更新至 $(0.6,1.8)$，损失由 $6$ 降为 $\answerresult{1.68}$。每步参数偏差分别乘以 $r_w=1-\tfrac43\eta$、$r_b=1-2\eta$：$\eta=1.2$ 时 $r_w=-0.6,r_b=-1.4$，$w$ 震荡收敛、$b$ 震荡发散；$\eta=0.9$ 时 $r_w=-0.2,r_b=-0.8$，两个方向都震荡收敛。
\item[\textbf{C-011.}] \phantomsection\label{solution:C-011}\label{solution:iterative-method-choice}\hyperref[exercise:iterative-method-choice]{\textbf{第一步：}}分别考虑寻找参数的方法与每次求梯度的数据量。\par\nobreak\textbf{解题过程：}本例直接求解更简便；它适合手算每一步，并用最优解核验迭代结果。一般多层网络的参数相互耦合，通常很难一次解出全部最优参数，而当前位置的梯度可以用于逐步更新。随机抽取单个样本或小批次，可以减少每次更新需要处理的数据量；估计的梯度受本次样本影响，全量训练损失可能暂时上升。Mini-batch 在每步计算量与梯度估计的波动之间折中。
\end{enumerate}
\Needspace{6\baselineskip}
\subsubsection*{第~\ref*{intro:07-functions-relu}~章 · 非线性：为什么多叠两层直线仍不够？\enspace{\mdseries\footnotesize（\hyperref[exercise:C-012]{题目见第~\pageref*{exercise:C-012}~页起}）}}
\begin{enumerate}[label={},leftmargin=4.8em,itemsep=4pt]
\item[\textbf{C-012.}] \phantomsection\label{solution:C-012}\label{solution:relu-determinism}\hyperref[exercise:relu-determinism]{\textbf{第一步：}}先计算 $z$，再分清参数 $w$ 与中间输出 $h$。\par\nobreak
\textbf{解题过程：}两次都是 $z=-1,h=0$；参数仍为 $w=-2,b=1$。ReLU 把负的中间输入映射为零，既不改写权重，也不引入随机性。固定输入和参数时，这一步结果相同。
\item[\textbf{C-013.}] \phantomsection\label{solution:C-013}\label{solution:relu-piecewise-slopes}\hyperref[exercise:relu-piecewise-slopes]{\textbf{第一步：}}先找两个折点 $x=1,3$，每一段只保留处于开启状态的单元。\par\nobreak
\textbf{解答：}$f(0)=0$，$f(2)=2$，$f(4)=2\times3-1=\answerresult{5}$；三段斜率依次为 $0$、$2$、$2-1=1$。每个 ReLU 在自己的折点处改变斜率，组合后形成两个折角。
\end{enumerate}
\Needspace{6\baselineskip}
\subsubsection*{第~\ref*{intro:17-backpropagation}~章 · 反向传播究竟在传什么？\enspace{\mdseries\footnotesize（\hyperref[exercise:C-014]{题目见第~\pageref*{exercise:C-014}~页起}）}}
\begin{enumerate}[label={},leftmargin=4.8em,itemsep=4pt]
\item[\textbf{C-014.}] \phantomsection\label{solution:C-014}\label{solution:local-chain-gradient}\hyperref[exercise:local-chain-gradient]{\textbf{第一步：}}每往前经过一步，就乘上那一步的局部导数。\par\nobreak
\textbf{解题过程：}$\partial\widehat{y}/\partial h=3$，所以 $\partial L/\partial h=(-2)\times3=\answerresult{-6}$；再乘 $\partial h/\partial w_1=x=2$，得 $\partial L/\partial w_1=(-6)\times2=\answerresult{-12}$。
\item[\textbf{C-015.}] \phantomsection\label{solution:C-015}\label{solution:scalar-chain-forward}\hyperref[exercise:scalar-chain-forward]{\textbf{第一步：}}先算 $h$ 和预测，再用链式法则求权重梯度。\par\nobreak
\textbf{解答：} $h=6$，$\widehat{y}=6$，$L=2$；$\partial L/\partial w_2=(\widehat{y}-y)h=(-2)\times6=-12$。
\item[\textbf{C-016.}] \phantomsection\label{solution:C-016}\label{solution:branch-gradient-sum}\hyperref[exercise:branch-gradient-sum]{\textbf{第一步：}}先算 $\widehat{y}$ 与损失对 $\widehat{y}$ 的梯度。\par\nobreak
\textbf{解题过程：}$\widehat{y}=5$，输出端梯度为 $5-4=1$。两条分支的贡献分别为 $1\times2=2$、$1\times3=3$，总梯度相加得 $\partial L/\partial h=\answerresult{5}$。若只取其中一条，就漏掉了 $h$ 对损失的另一条影响路径。
\item[\textbf{C-017.}] \phantomsection\label{solution:C-017}\label{solution:relu-blocked-gradient}\hyperref[exercise:relu-blocked-gradient]{\textbf{第一步：}}先判断 $z$ 落在 ReLU 的哪一段。\par\nobreak
\textbf{解题过程：}$z=-2,h=0,\widehat{y}=0,L=1/2$。虽然 $\partial L/\partial h=(0-1)\times3=-3$，但 $z<0$ 时 $\partial h/\partial z=0$，所以 $\partial L/\partial w=(-3)\times0\times2=\answerresult{0}$。损失非零，仍可能被中间某一步的零导数截断。
\item[\textbf{C-018.}] \phantomsection\label{solution:C-018}\label{solution:branch-gradient-reuse}\hyperref[exercise:branch-gradient-reuse]{\textbf{第一步：}}先把到达同一节点的各分支梯度汇总。\par\nobreak\textbf{解题过程：}同一节点收到来自多个输出分支的梯度后求和，再向前传播；可复用保存的中间值与局部导数。
\end{enumerate}
\Needspace{6\baselineskip}
\subsubsection*{第~\ref*{intro:18-generalization}~章 · 训练得很好，为什么新题反而做错？\enspace{\mdseries\footnotesize（\hyperref[exercise:C-019]{题目见第~\pageref*{exercise:C-019}~页起}）}}
\begin{enumerate}[label={},leftmargin=4.8em,itemsep=4pt]
\item[\textbf{C-019.}] \phantomsection\label{solution:C-019}\label{solution:validation-checkpoint-choice}\hyperref[exercise:validation-checkpoint-choice]{\textbf{第一步：}}先只比较验证损失，不用训练损失决定胜者。\par\nobreak
\textbf{解题过程：}应选验证损失最小的 $\answerresult{\text{C}}$，而不是训练损失最小的 B。用最终测试结果参与选择后，测试数据已影响选模；应另留独立测试数据做最终评价。
\item[\textbf{C-020.}] \phantomsection\label{solution:C-020}\label{solution:zero-loss-generalization}\hyperref[exercise:zero-loss-generalization]{\textbf{第一步：}}先分别看训练样本与独立样本。\par\nobreak
\textbf{解题过程：}$y=4$ 时 $\tfrac12(4-4)^2=0$，只说明这一份样本被拟合；仍要用独立验证数据检验，并排除改标签或数据泄漏造成的假象。
\item[\textbf{C-021.}] \phantomsection\label{solution:C-021}\label{solution:validation-overfitting}\hyperref[exercise:validation-overfitting]{\textbf{第一步：}}对比训练和验证曲线的变化方向。\par\nobreak
\textbf{解答：} 可能过拟合；可考虑更多合适数据、数据增强、L2 惩罚、Dropout、提前停止等，需用独立验证数据评估。
\item[\textbf{C-022.}] \phantomsection\label{solution:C-022}\label{solution:small-data-validation}\hyperref[exercise:small-data-validation]{\textbf{第一步：}}先把最终测试集保留在选择过程之外。\par\nobreak\textbf{解题过程：}在训练集内部采用交叉验证或重复验证切分做模型选择，最终测试集独立保留；代价是多次训练成本。
\end{enumerate}
\Needspace{6\baselineskip}
\subsubsection*{第~\ref*{intro:01-vectors}~章 · 向量与点积：怎样把多个特征变成一个匹配分数？\enspace{\mdseries\footnotesize（\hyperref[exercise:C-023]{题目见第~\pageref*{exercise:C-023}~页起}）}}
\begin{enumerate}[label={},leftmargin=4.8em,itemsep=4pt]
\item[\textbf{C-023.}] \phantomsection\label{solution:C-023}\label{solution:vector-coordinate-operations}\hyperref[exercise:vector-coordinate-operations]{\textbf{第一步：}}加法与数乘都逐个坐标计算；长度用勾股定理。\par\nobreak
\textbf{解答：}$\vect a^\top+\vect b^\top=(4,2)$，$2\vect a^\top=(6,8)$，$\|\vect a\|=\sqrt{3^2+4^2}=\answerresult{5}$。交换坐标后数值集合不变，但每个位置对应的特征变了，一般不再表示同一对象。
\item[\textbf{C-024.}] \phantomsection\label{solution:C-024}\label{solution:unit-vector-dot-comparison}\hyperref[exercise:unit-vector-dot-comparison]{\textbf{第一步：}}按同一特征位置相乘，再把两项相加。\par\nobreak
\textbf{解题过程：}$\vect a^\top\vect b=1\times0+0\times1=\answerresult{0}$，$\vect a^\top\vect c=1\times0.6+0\times0.8=\answerresult{0.6}$，所以 C 的分数更高。点积计算本身不要求单位长度；本题三条向量恰好长度均为 $1$，因此分数也可直接比较方向的接近程度。
\item[\textbf{C-025.}] \phantomsection\label{solution:C-025}\label{solution:article-query-ranking}\hyperref[exercise:article-query-ranking]{\textbf{第一步：}}用新的需求分别与三篇文章做点积。\par\nobreak\textbf{解题过程：}$\vect q^\top\vect a=0.8\times0.8+0.6\times0.6=\answerresult{1}$，$\vect q^\top\vect b=\answerresult{0.8}$，$\vect q^\top\vect c=\answerresult{0.6}$。排名为 A、B、C；技术项的需求权重高于生活项，因而纯技术的 B 超过纯生活的 C。
\item[\textbf{C-026.}] \phantomsection\label{solution:C-026}\label{solution:cosine-rescaling}\hyperref[exercise:cosine-rescaling]{\textbf{第一步：}}先算原始点积，再除以两个长度的乘积。\par\nobreak\textbf{解题过程：}点积 $3\times4+4\times3=\answerresult{24}$，两个长度均为 $\sqrt{3^2+4^2}=\answerresult{5}$；余弦相似度为 $24/(5\times5)=\answerresult{0.96}$。两条向量各放大 $5$ 倍，原始点积放大 $25$ 倍，方向与夹角不变；分母也放大 $25$ 倍，所以余弦相似度保持不变。
\item[\textbf{C-027.}] \phantomsection\label{solution:C-027}\label{solution:cosine-length-control}\hyperref[exercise:cosine-length-control]{\textbf{第一步：}}先把每条向量除以自己的长度。\par\nobreak\textbf{解题过程：}余弦相似度除以两个范数，排除了非零向量整体放大的影响；原始点积则会随长度变化。
\end{enumerate}
\Needspace{6\baselineskip}
\subsubsection*{第~\ref*{intro:03-matrix-multiplication}~章 · 矩阵乘法到底在算什么？\enspace{\mdseries\footnotesize（\hyperref[exercise:C-028]{题目见第~\pageref*{exercise:C-028}~页起}）}}
\begin{enumerate}[label={},leftmargin=4.8em,itemsep=4pt]
\item[\textbf{C-028.}] \phantomsection\label{solution:C-028}\label{solution:matrix-single-entry}\hyperref[exercise:matrix-single-entry]{\textbf{第一步：}}结果的第 $i$ 行第 $j$ 列，来自左边第 $i$ 行乘右边第 $j$ 列。\par\nobreak
\textbf{解题过程：}$AB:2\times3$。$(AB)_{12}=1\times(-1)+2\times3=\answerresult{5}$；$(AB)_{23}=3\times1+0\times0=\answerresult{3}$。它们分别用了 $A$ 的第 $1$ 行与 $B$ 的第 $2$ 列、$A$ 的第 $2$ 行与 $B$ 的第 $3$ 列。
\item[\textbf{C-029.}] \phantomsection\label{solution:C-029}\label{solution:matrix-batch-calculation}\hyperref[exercise:matrix-batch-calculation]{\textbf{第一步：}}先让第一行输入乘 $W$ 的两列，再处理第二行。\par\nobreak
\textbf{解题过程：}$Y=\answerresult{\begin{bmatrix}-2&2\\5&4\end{bmatrix}}$，两行结果分别为 $(1-3,2)=(-2,2)$ 与 $(2+3,4)=(5,4)$。每行独立使用同一个 $W$；只改第一行输入，第二行仍为 $(5,4)$。
\item[\textbf{C-030.}] \phantomsection\label{solution:C-030}\label{solution:matrix-shape-correction}\hyperref[exercise:matrix-shape-correction]{\textbf{第一步：}}相乘前先比较左矩阵的列数与右矩阵的行数；不相等时考虑转置哪一个。\par\nobreak
\textbf{解题过程：}$AB$ 的内侧维度为 $2$ 和 $5$，不能相乘。改为 $AB^\top:3\times5$，再乘 $C:5\times4$，得 $AB^\top C:\answerresult{3\times4}$；其转置 $(AB^\top C)^\top=C^\top BA^\top$ 为 $4\times3$。
\item[\textbf{C-031.}] \phantomsection\label{solution:C-031}\label{solution:batch-matrix-sharing}\hyperref[exercise:batch-matrix-sharing]{\textbf{第一步：}}先算一行的输出宽度，再数行数。\par\nobreak\textbf{解题过程：}每行 $1\times2$ 乘 $2\times3$ 得 $1\times3$，五行堆叠是 $5\times3$；整批一次 $XW$ 也得 $5\times3$。
\item[\textbf{C-032.}] \phantomsection\label{solution:C-032}\label{solution:gram-asymmetry}\hyperref[exercise:gram-asymmetry]{\textbf{第一步：}}$G$ 的第 $i$ 行第 $j$ 列是 $X$ 第 $i$ 行与第 $j$ 行的点积；再用 $(AB)^\top=B^\top A^\top$ 检查。\par\nobreak
\textbf{解题过程：}$G=\begin{bmatrix}5&2&-1\\2&1&-1\\-1&-1&2\end{bmatrix}$；由 $G^\top=(X^\top)^\top X^\top=XX^\top$，它必然对称。$XW=\begin{bmatrix}1&3\\0&1\\1&0\end{bmatrix}$，$S=\begin{bmatrix}7&3&-2\\2&1&-1\\1&0&1\end{bmatrix}$，例如 $S_{13}=-2\neq S_{31}=1$；因为 $S^\top=XW^\top X^\top$，只有 $W^\top=W$ 时才等于 $S$。
\end{enumerate}
\Needspace{6\baselineskip}
\subsubsection*{第~\ref*{intro:10-mlp}~章 · MLP：把矩阵与激活函数组成两层网络\enspace{\mdseries\footnotesize（\hyperref[exercise:C-033]{题目见第~\pageref*{exercise:C-033}~页起}）}}
\begin{enumerate}[label={},leftmargin=4.8em,itemsep=4pt]
\item[\textbf{C-033.}] \phantomsection\label{solution:C-033}\label{solution:mlp-activation-forward}\hyperref[exercise:mlp-activation-forward]{\textbf{第一步：}}先算两个隐藏单元的加权和，记得分别加偏置。\par\nobreak
\textbf{解题过程：}$\vect z^\top=(2-2+1,-1-2)=(1,-3)$，$\vect h^\top=\answerresult{(1,0)}$，因此 $\widehat{y}=1\times2+0\times(-1)+0.5=\answerresult{2.5}$。两个隐藏分量只保留了正的第一项。
\item[\textbf{C-034.}] \phantomsection\label{solution:C-034}\label{solution:mlp-parameter-count}\hyperref[exercise:mlp-parameter-count]{\textbf{第一步：}}按输入宽度乘输出宽度数权重，每个输出单元再配一个偏置。\par\nobreak
\textbf{解题过程：}$W_1:2\times3,\vect b_1^\top:1\times3,W_2:3\times1,b_2:1\times1$。总参数为 $2\times3+3+3\times1+1=\answerresult{13}$。五个样本共用参数，参数量仍为 $13$；输入 $5\times2$ 经隐藏层 $5\times3$ 后得到输出 $\answerresult{5\times1}$。
\item[\textbf{C-035.}] \phantomsection\label{solution:C-035}\label{solution:mlp-negative-activation}\hyperref[exercise:mlp-negative-activation]{\textbf{第一步：}}先算上投影，找出会被 ReLU 截断的负分量。\par\nobreak
\textbf{解答：}上投影是 $(2,1,-1)$，ReLU 后为 $(2,1,0)$，下投影得 $\answerresult{(4,1)}$；去掉 ReLU 则 $(2,1,-1)$ 直接下投影，得 $(3,0)$。差别全部来自被置零的第三个分量。
\item[\textbf{C-036.}] \phantomsection\label{solution:C-036}\label{solution:mlp-linear-composition}\hyperref[exercise:mlp-linear-composition]{\textbf{第一步：}}先试把相邻的两个矩阵合并。\par\nobreak\textbf{解题过程：}无激活时可合并为 $x(W_1W_2)$；含 ReLU 时计算分支依赖输入，能表达固定线性矩阵无法表达的函数。
\end{enumerate}
\Needspace{6\baselineskip}
\subsubsection*{第~\ref*{intro:20a-tokenizer}~章 · Token 与 ID：文字怎样进入数字世界？\enspace{\mdseries\footnotesize（\hyperref[exercise:C-037]{题目见第~\pageref*{exercise:C-037}~页起}）}}
\begin{enumerate}[label={},leftmargin=4.8em,itemsep=4pt]
\item[\textbf{C-037.}] \phantomsection\label{solution:C-037}\label{solution:bpe-pair-merge}\hyperref[exercise:bpe-pair-merge]{\textbf{第一步：}}独立序列内部统计相邻对，比较次数。\par\nobreak
\textbf{解答：} 两段内相邻对的总计数为
$(a,a):1$，$(a,b):2$；第一轮合并 $a+b\to ab$，
得到 \texttt{a ab} 与 \texttt{ab}。
每段单独切分与合并，不跨段统计相邻对。
\item[\textbf{C-038.}] \phantomsection\label{solution:C-038}\label{solution:tokenizer-attention-cost}\hyperref[exercise:tokenizer-attention-cost]{\textbf{第一步：}}先把长度代入 $T\times T$。\par\nobreak\textbf{解题过程：}分数矩阵由 $T^2$ 增到 $(2T)^2=\answerresult{4T^2}$，元素数四倍；端到端耗时还受实现、投影、硬件和序列生成方式等影响。
\end{enumerate}
\Needspace{6\baselineskip}
\subsubsection*{第~\ref*{intro:20b-task-contract}~章 · 把语言练习变成一道可计算的预测题\enspace{\mdseries\footnotesize（\hyperref[exercise:C-039]{题目见第~\pageref*{exercise:C-039}~页起}）}}
\begin{enumerate}[label={},leftmargin=4.8em,itemsep=4pt]
\item[\textbf{C-039.}] \phantomsection\label{solution:C-039}\label{solution:shifted-token-targets}\hyperref[exercise:shifted-token-targets]{\textbf{第一步：}}训练时输入序列和目标序列错开一个位置。\par\nobreak
\textbf{解答：} 输入 $(2,4,1)$，目标 $(4,1,8)$，且所有 ID 均在词表范围内。
\item[\textbf{C-040.}] \phantomsection\label{solution:C-040}\label{solution:causal-target-visibility}\hyperref[exercise:causal-target-visibility]{\textbf{第一步：}}从当前位置向前列出允许读取的位置，再找到目标与预测汇合的步骤。\par\nobreak
\textbf{解答：} 三个位置分别允许读取位置 $\answerresult{\{1\}}$、$\answerresult{\{1,2\}}$、$\answerresult{\{1,2,3\}}$。
第三位置的目标 ID $3$ 指出该行应预测“。”；评分步骤比较该行预测与这个答案，计算预测损失。
\item[\textbf{C-041.}] \phantomsection\label{solution:C-041}\label{solution:training-generation-data}\hyperref[exercise:training-generation-data]{\textbf{第一步：}}分别追踪训练目标和生成时新输入的来源。\par\nobreak\textbf{解题过程：}训练目标来自已有文本：输入 $[2,4,5]$，对应目标是 $[4,5,3]$，即“花、开了、。”。
生成时只有给定的“春天”，模型先输出下一 Token 的概率，再按选取规则得到一个 Token 并追加到输入；若选出“花”，下一轮输入就是“春天 花”。
训练目标从文本直接取得，生成时追加的 Token 来自模型的预测与选取。
\end{enumerate}
\Needspace{6\baselineskip}
\subsubsection*{第~\ref*{intro:20c-classification-loss}~章 · 分类损失：预测错了哪一类，怎样调整分数？\enspace{\mdseries\footnotesize（\hyperref[exercise:C-042]{题目见第~\pageref*{exercise:C-042}~页起}）}}
\begin{enumerate}[label={},leftmargin=4.8em,itemsep=4pt]
\item[\textbf{C-042.}] \phantomsection\label{solution:C-042}\label{solution:single-token-cross-entropy}\hyperref[exercise:single-token-cross-entropy]{\textbf{第一步：}}利用 $-\ln p=\ln(1/p)$，先把概率写成倒数。\par\nobreak
\textbf{解答：} $-\ln(0.25)=\ln4\approx1.386$。
\item[\textbf{C-043.}] \phantomsection\label{solution:C-043}\label{solution:softmax-logit-gradient}\hyperref[exercise:softmax-logit-gradient]{\textbf{第一步：}}先把类别 1 写成与 $p$ 同形的 One-hot 向量。\par\nobreak
\textbf{解答：}$y=(1,0,0)$，$p-y=\answerresult{(-0.8,0.5,0.3)}$。沿负梯度更新，$z_1$ 升高，$z_2,z_3$ 降低；错误类中概率越大的 $z_2$ 被压低得越多。
\item[\textbf{C-044.}] \phantomsection\label{solution:C-044}\label{solution:probability-loss-ratio}\hyperref[exercise:probability-loss-ratio]{\textbf{第一步：}}先分别取负对数，再用 $\ln a-\ln b=\ln(a/b)$。\par\nobreak
\textbf{解答：}$-\ln0.2\approx1.609$，$-\ln0.8\approx0.223$，差约 $\answerresult{1.386}=\ln(0.8/0.2)=\ln4$：正确类概率提高到四倍，惊讶度减少 $\ln4$。
\item[\textbf{C-045.}] \phantomsection\label{solution:C-045}\label{solution:one-hot-gradient-target}\hyperref[exercise:one-hot-gradient-target]{\textbf{第一步：}}先把类别标签写成与 $p$ 同形的向量。\par\nobreak\textbf{解题过程：}$y$ 是与概率同形的 one-hot 目标向量，正确类分量为 $1$、其他为 $0$；整数 ID 只用于索引。
\end{enumerate}
\Needspace{6\baselineskip}
\subsubsection*{第~\ref*{intro:20-embeddings}~章 · Embedding：让 Token ID 变成可训练的向量\enspace{\mdseries\footnotesize（\hyperref[exercise:C-046]{题目见第~\pageref*{exercise:C-046}~页起}）}}
\begin{enumerate}[label={},leftmargin=4.8em,itemsep=4pt]
\item[\textbf{C-046.}] \phantomsection\label{solution:C-046}\label{solution:embedding-input-change}\hyperref[exercise:embedding-input-change]{\textbf{第一步：}}按新输入的三个 ID 依次取表行。\par\nobreak
\textbf{解题过程：}新的查表结果为
\[\begin{bmatrix}0.15&-0.2&1.25\\-0.3&1.1&0.2\\0.15&-0.2&1.25\end{bmatrix}.\]
只有第一行查表结果改变；第 1、3 行因为读取同一个 ID 而相同。这里只更换输入，没有优化器更新，Embedding 表中的参数保持不变。
\item[\textbf{C-047.}] \phantomsection\label{solution:C-047}\label{solution:embedding-shape}\hyperref[exercise:embedding-shape]{\textbf{第一步：}}词表项数确定行数，特征宽度确定列数。\par\nobreak
\textbf{解答：} $E:4\times3$；每个 Token 嵌入为 3 维。
\item[\textbf{C-048.}] \phantomsection\label{solution:C-048}\label{solution:embedding-update}\hyperref[exercise:embedding-update]{\textbf{第一步：}}把给定坐标、梯度和学习率代入“参数减去学习率乘梯度”。\par\nobreak
\textbf{解题过程：}$-0.3-0.04\times0.257642=\answerresult{-0.31030568}\approx-0.31031$。更换输入 ID 只改变读取哪一行；没有参数更新时，$E[4]$ 保持原值。
\end{enumerate}
\Needspace{6\baselineskip}
\subsubsection*{第~\ref*{intro:21a-context-gap}~章 · 前文怎样传到当前词？从逐行处理到 RNN\enspace{\mdseries\footnotesize（\hyperref[exercise:C-049]{题目见第~\pageref*{exercise:C-049}~页起}）}}
\begin{enumerate}[label={},leftmargin=4.8em,itemsep=4pt]
\item[\textbf{C-049.}] \phantomsection\label{solution:C-049}\label{solution:history-model-comparison}\hyperref[exercise:history-model-comparison]{\textbf{第一步：}}分别画出固定窗口包含的输入，以及 RNN 的状态传递箭头。\par\nobreak\textbf{解题过程：}固定窗口把最近若干 Token 的向量拼成输入，窗口外的信息不能直接进入。RNN 每步接收当前输入与上一状态，把部分历史压进固定宽度的摘要；本章的递推不会向当前步骤提供所有历史状态供单独读取。
\item[\textbf{C-050.}] \phantomsection\label{solution:C-050}\label{solution:window-input-width}\hyperref[exercise:window-input-width]{\textbf{第一步：}}拼接把各行首尾相接，宽度等于窗口长度乘每行维度。\par\nobreak
\textbf{解答：}宽度分别为 $2\times3=6$ 与 $8\times3=\answerresult{24}$。$n=2$ 只覆盖前一位置与当前位置，看不到距离 $4$ 的线索；$n=8$ 能看到。但输入宽度和首层权重随 $n$ 线性增长，窗口外的信息仍会丢失。
\item[\textbf{C-051.}] \phantomsection\label{solution:C-051}\label{solution:rnn-decay}\hyperref[exercise:rnn-decay]{\textbf{第一步：}}只有第一个输入非零，之后每一步状态只乘 $0.8$。\par\nobreak
\textbf{解答：}$h_1=1,h_2=0.8,h_3=0.64,h_4=0.512$，一般 $h_t=0.8^{t-1}$。$0.8^{10}\approx0.107$，$0.8^{11}\approx0.086$，故从 $t=\answerresult{12}$ 起小于 $0.1$。
\item[\textbf{C-052.}] \phantomsection\label{solution:C-052}\label{solution:rnn-state-scaling}\hyperref[exercise:rnn-state-scaling]{\textbf{第一步：}}把递推式 $h_t=f(x_t,h_{t-1})$ 的输入与输出形状写出。\par\nobreak\textbf{解题过程：}单个状态仍为 $1\times d$，更新次数从 $T$ 次增至 $2T$ 次。最后一步直接接收 $h_{2T-1}$；较早的信息须经过中间状态持续保留在摘要里。这要求固定宽度的状态能保住任务所需的历史线索。
\end{enumerate}
\Needspace{6\baselineskip}
\subsubsection*{第~\ref*{intro:21c-naive-attention}~章 · Attention：从直接读取到 Q/K/V 分工\enspace{\mdseries\footnotesize（\hyperref[exercise:C-053]{题目见第~\pageref*{exercise:C-053}~页起}）}}
\begin{enumerate}[label={},leftmargin=4.8em,itemsep=4pt]
\item[\textbf{C-053.}] \phantomsection\label{solution:C-053}\label{solution:weighted-value-read}\hyperref[exercise:weighted-value-read]{\textbf{第一步：}}先把两条 Value 按行排成矩阵，输出是两行的加权和。\par\nobreak
\textbf{解题过程：}$V=\begin{bmatrix}2&1\\0&3\end{bmatrix}$，$aV=0.8(2,1)+0.2(0,3)=\answerresult{(1.6,1.4)}$；矩阵相乘逐列得到同样两项。权重和为 $1$。换成 $(0.25,0.75)$ 后得到 $(0.5,0.25+2.25)=\answerresult{(0.5,2.5)}$。
\item[\textbf{C-054.}] \phantomsection\label{solution:C-054}\label{solution:unnormalized-value-read}\hyperref[exercise:unnormalized-value-read]{\textbf{第一步：}}先取匹配权重，再按行加权求和 Value。\par\nobreak
\textbf{解答：} $aV=(0,6)$；这是两行 Value 按权重 $0$ 与 $2$ 求和。
\item[\textbf{C-055.}] \phantomsection\label{solution:C-055}\label{solution:self-dot-length}\hyperref[exercise:self-dot-length]{\textbf{第一步：}}分别计算 $x_1\cdot x_3$、$x_3\cdot x_1$ 与 $x_3\cdot x_3$。\par\nobreak\textbf{解题过程：}$S_{13}=\answerresult{2}$、$S_{31}=\answerresult{2}$、$S_{33}=\answerresult{8}$；$XX^\top$ 仍然对称。方向相同的 $x_3$ 长度增大一倍，与 $x_1$ 的点积增大一倍，与自身的点积增大四倍，说明分数也受向量长度影响。
\item[\textbf{C-056.}] \phantomsection\label{solution:C-056}\label{solution:query-projection-choice}\hyperref[exercise:query-projection-choice]{\textbf{第一步：}}查询向量可以来自可学习投影。\par\nobreak
\textbf{解答：} 简化模型直接使用输入作为查询；标准 Transformer 通过可学习投影 $\vect q^\top=\vect x^\top W_Q$ 得到查询。
\item[\textbf{C-057.}] \phantomsection\label{solution:C-057}\label{solution:key-value-projections}\hyperref[exercise:key-value-projections]{\textbf{第一步：}}先比较匹配特征与内容特征是否必须相同。\par\nobreak\textbf{解题过程：}把匹配和内容绑定限制灵活度；独立 $W_Q,W_K,W_V$ 能分别学习匹配视角和提供的内容。
\end{enumerate}
\Needspace{6\baselineskip}
\subsubsection*{第~\ref*{intro:11-softmax-memory}~章 · 从单 Query 到批量 Self-Attention\enspace{\mdseries\footnotesize（\hyperref[exercise:C-058]{题目见第~\pageref*{exercise:C-058}~页起}）}}
\begin{enumerate}[label={},leftmargin=4.8em,itemsep=4pt]
\item[\textbf{C-058.}] \phantomsection\label{solution:C-058}\label{solution:attention-loss-transfer}\hyperref[exercise:attention-loss-transfer]{\textbf{第一步：}}按点积、Softmax、读出、输出头、交叉熵的顺序算。\par\nobreak
\textbf{解答：} $qK^\top=(1,1,2)$；Softmax 为
$(e,e,e^2)/(2e+e^2)\approx(0.212,0.212,0.576)$，
读取 $r\approx(0.788058,0.788058)$。
按照题设矩阵右乘，logits 约为
$z=(1.576117,0.788058,0)$，再做 Softmax，
第二类概率约 $p_2=0.273684$，所以
$\answerresult{L=-\ln p_2\approx1.295781}$。
\textbf{这次预测真正沿着读取向量和输出头得出}；
若不给出输出头，Attention 权重本身不足以唯一决定词表概率。
\item[\textbf{C-059.}] \phantomsection\label{solution:C-059}\label{solution:attention-shapes}\hyperref[exercise:attention-shapes]{\textbf{第一步：}}依次核对 $QK^\top$、权重以及权重乘 $V$ 的形状。\par\nobreak
\textbf{解答：} $QK^\top:3\times5$；权重仍为 $3\times5$；最终输出 $3\times6$。
\item[\textbf{C-060.}] \phantomsection\label{solution:C-060}\label{solution:batch-softmax-recompute}\hyperref[exercise:batch-softmax-recompute]{\textbf{第一步：}}每行先做指数归一化，再右乘同一 $V$。\par\nobreak
\textbf{解答：} 每行指数分数分别为 $(2,1,1)$、$(1,3,1)$、$(1,1,4)$，
故 $A'=\boxed{\begin{bmatrix}
1/2&1/4&1/4\\
1/5&3/5&1/5\\
1/6&1/6&2/3
\end{bmatrix}}$。
与题设 $V:3\times2$ 相乘，
$Y'=\boxed{\begin{bmatrix}
3/4&1/2\\
2/5&4/5\\
5/6&5/6
\end{bmatrix}}$，形状为 $3\times2$。
每行权重和为 $1$；本题与正文的分数矩阵不同，不能直接抄原来的 $A$。
\item[\textbf{C-061.}] \phantomsection\label{solution:C-061}\label{solution:cross-attention-shapes}\hyperref[exercise:cross-attention-shapes]{\textbf{第一步：}}先得分矩阵，再看它右乘 $V$ 后的形状。\par\nobreak
\textbf{解答：} $QK^\top:2\times4$，输出 $2\times5$；Softmax 对每个 Query 的\textbf{一行} Key 分数分别归一化。
\item[\textbf{C-062.}] \phantomsection\label{solution:C-062}\label{solution:score-value-confusion}\hyperref[exercise:score-value-confusion]{\textbf{第一步：}}先区分分数矩阵与读出矩阵。\par\nobreak\textbf{解题过程：}$QK^\top$ 只生成分数而非内容，须先用 Softmax 得权重，再右乘 Value 才得到读出表示。
\end{enumerate}
\Needspace{6\baselineskip}
\subsubsection*{第~\ref*{intro:22a-position-multihead}~章 · 位置编码、缩放与多头注意力的手算\enspace{\mdseries\footnotesize（\hyperref[exercise:C-063]{题目见第~\pageref*{exercise:C-063}~页起}）}}
\begin{enumerate}[label={},leftmargin=4.8em,itemsep=4pt]
\item[\textbf{C-063.}] \phantomsection\label{solution:C-063}\label{solution:head-concatenation}\hyperref[exercise:head-concatenation]{\textbf{第一步：}}多头结果按列拼接，再确定输出矩阵的输入宽度。\par\nobreak
\textbf{解答：} 拼接后 $3\times4$；若按行向量约定，$W_O:4\times4$。
\item[\textbf{C-064.}] \phantomsection\label{solution:C-064}\label{solution:head-projection-diversity}\hyperref[exercise:head-projection-diversity]{\textbf{第一步：}}先比较两个头是否拥有独立参数。\par\nobreak\textbf{解题过程：}每个头可以分别训练 Q/K/V 投影，读取不同子空间；复制相同头不创造独立表示能力。
\end{enumerate}
\Needspace{6\baselineskip}
\subsubsection*{第~\ref*{intro:22-transformer-architecture}~章 · 把 Attention 组装起来：Transformer 不只是一个公式\enspace{\mdseries\footnotesize（\hyperref[exercise:C-065]{题目见第~\pageref*{exercise:C-065}~页起}）}}
\begin{enumerate}[label={},leftmargin=4.8em,itemsep=4pt]
\item[\textbf{C-065.}] \phantomsection\label{solution:C-065}\label{solution:causal-position-visibility}\hyperref[exercise:causal-position-visibility]{\textbf{第一步：}}因果掩码只允许当前位置读取自己及此前位置。\par\nobreak
\textbf{解答：} 第 2 个位置能看到位置 $1,2$；第 1 个位置不能看到位置 $4$。
\item[\textbf{C-066.}] \phantomsection\label{solution:C-066}\label{solution:block-branch-roles}\hyperref[exercise:block-branch-roles]{\textbf{第一步：}}分别标出跨位置、逐位置与保留通路。\par\nobreak\textbf{解题过程：}Attention 进行跨位置读取，FFN 每位置改写特征；残差帮助保留信息与梯度通路。
\item[\textbf{C-067.}] \phantomsection\label{solution:C-067}\label{solution:token-to-loss-pipeline}\hyperref[exercise:token-to-loss-pipeline]{\textbf{第一步：}}先从 ID 查表画到 logits 和损失。\par\nobreak\textbf{解题过程：}Token ID 查表加位置得到隐藏矩阵，Block 输出上下文化隐藏状态，输出头变 logits，Softmax 变概率；训练把概率与目标比较得到损失，再反传更新参数。
\end{enumerate}
\Needspace{6\baselineskip}
\subsubsection*{第~\ref*{intro:10b-ffn-output-head}~章 · 前馈网络与输出头：两种不同的特征变换\enspace{\mdseries\footnotesize（\hyperref[exercise:C-068]{题目见第~\pageref*{exercise:C-068}~页起}）}}
\begin{enumerate}[label={},leftmargin=4.8em,itemsep=4pt]
\item[\textbf{C-068.}] \phantomsection\label{solution:C-068}\label{solution:ffn-output-shapes}\hyperref[exercise:ffn-output-shapes]{\textbf{第一步：}}FFN 先扩宽再恢复 $d_{\rm model}$；只有输出头才出现词表轴。\par\nobreak
\textbf{解答：}$W_{\rm up}:8\times32$，$W_{\rm down}:32\times8$，FFN 输出为 $5\times8$；$W_{\rm out}:8\times50$，logits 为 $\answerresult{5\times50}$。
\item[\textbf{C-069.}] \phantomsection\label{solution:C-069}\label{solution:ffn-weight-count}\hyperref[exercise:ffn-weight-count]{\textbf{第一步：}}分别数两张 FFN 权重表的元素。\par\nobreak\textbf{解题过程：}$6\cdot18+18\cdot6=\answerresult{216}$ 个权重；FFN 宽度与词表大小不同，词表输出头另计。
\end{enumerate}
\Needspace{6\baselineskip}
\subsubsection*{第~\ref*{intro:22b-token-to-output}~章 · 三维贯通实验：从 Token ID 到三个位置的输出概率\enspace{\mdseries\footnotesize（\hyperref[exercise:C-070]{题目见第~\pageref*{exercise:C-070}~页起}）}}
\begin{enumerate}[label={},leftmargin=4.8em,itemsep=4pt]
\item[\textbf{C-070.}] \phantomsection\label{solution:C-070}\label{solution:causal-input-change}\hyperref[exercise:causal-input-change]{\textbf{第一步：}}按因果 Mask 写出每一行能读取的输入位置。\par\nobreak\textbf{解题过程：}（a）第 $i$ 行只读取前 $i$ 个位置，LayerNorm 与 FFN 又逐行计算，所以第 1、2 行完全不变；第 3 行的输入与读取内容都变了，概率随之改变（复算得第 3 行最高概率改为 ID 4，约 $0.7212$）。（b）$X,H:\answerresult{4\times3}$，$A:\answerresult{4\times4}$，$Z,\Pi:\answerresult{4\times5}$；但 $P_{\rm pos}$ 只有 $3$ 行，缺少第 $4$ 个位置向量，主例的上下文窗口只能容纳 $3$ 个 Token。
\item[\textbf{C-071.}] \phantomsection\label{solution:C-071}\label{solution:output-head-only-training}\hyperref[exercise:output-head-only-training]{\textbf{第一步：}}只训练输出头时，三行隐藏表示 $H$ 保持初值，变化的只有最后的线性读出表。\par\nobreak
\textbf{解题过程：}实测第 20 步平均交叉熵为 $\answerresult{0.816707}$，贪心预测 $[4,4,3]$，位置 2 仍错（目标 ID 5）。完整训练同一步为 $0.483173$、$[4,5,3]$，第 9 步起三位置全对。固定的隐藏行限制了只调输出头在 20 步内能达到的结果；改变 $H$ 需要前面各层也接收梯度。
\end{enumerate}
\Needspace{6\baselineskip}
\subsubsection*{第~\ref*{intro:22c-language-training}~章 · 三个位置一起训练：从交叉熵到共享参数更新\enspace{\mdseries\footnotesize（\hyperref[exercise:C-072]{题目见第~\pageref*{exercise:C-072}~页起}）}}
\begin{enumerate}[label={},leftmargin=4.8em,itemsep=4pt]
\item[\textbf{C-072.}] \phantomsection\label{solution:C-072}\label{solution:context-conflicting-labels}\hyperref[exercise:context-conflicting-labels]{\textbf{第一步：}}比较两条样本送入输出头的特征是否仍不同。\par\nobreak\textbf{解题过程：}两条前文相同后，当前词、位置和读取结果全部相同，两个版本都只能给出同一分布；等权相反二类标签的最优平均损失为 $\ln2$。容量不能补回输入中没有的区分信息。
\item[\textbf{C-073.}] \phantomsection\label{solution:C-073}\label{solution:rational-mini-gradient}\hyperref[exercise:rational-mini-gradient]{\textbf{第一步：}}从二类交叉熵的 $g_z=(1/2,-1/2)$ 出发。\par\nobreak\textbf{解题过程：}$g_h=(1,-1/3)$；FFN 两个单位投影与正区间 ReLU 传回同一梯度，因此 $g_r=2g_h=(2,-2/3)$；Softmax 分数梯度 $(2/3,-2/3)$ 乘 Query $q=1$，得到 $\nabla_{W_K}L=(2/3,-2/3)^\top$。完整参数导数见第~\pageref*{app:rational-mini}~页。
\item[\textbf{C-074.}] \phantomsection\label{solution:C-074}\label{solution:simplified-block-gradient}\hyperref[exercise:simplified-block-gradient]{\textbf{第一步：}}先求 $f^\top g_h$，再沿残差路径反传。\par\nobreak
\textbf{解答：} 本题已关闭 Norm、固定 $W_{\rm up}=I_2$；$r=(0.330,1.670)$ 两项为正，故 $f=\relu(r)=r$，$h=r+fW_{\rm down}$，
$\partial L/\partial W_{\rm down}=f^\top g_h
\approx\begin{bmatrix}0.03206&0.23930\\0.16226&1.21102\end{bmatrix}$。
上游梯度 $g_r=g_h(I_2+0.5I_2)=1.5g_h$，
故 $\partial L/\partial W_V=a^\top g_r
\approx\begin{bmatrix}0.04809&0.35895\\0.09765&0.72879\end{bmatrix}$。
\item[\textbf{C-075.}] \phantomsection\label{solution:C-075}\label{solution:correct-probability-loss}\hyperref[exercise:correct-probability-loss]{\textbf{第一步：}}比较同一正确类两次的 $-\ln p$。\par\nobreak\textbf{解题过程：}$-\ln0.07<-\ln0.05$，正确类概率上升意味着该位置的交叉熵下降，即使贪心类别暂时未改正。
\item[\textbf{C-076.}] \phantomsection\label{solution:C-076}\label{solution:attention-attribution}\hyperref[exercise:attention-attribution]{\textbf{第一步：}}只改变一组参数是否可训练，其余设置保持不变。\par\nobreak\textbf{解题过程：}加 \texttt{-{}-freeze Wq,Wk} 冻结 $W_Q,W_K$，$A_{3,2}$ 仍升至 $0.324789$；加 \texttt{-{}-only Wq,Wk} 只训练它们，$A_{3,2}$ 反降至 $0.223962$。变化主要来自 $E$、$P_{\rm pos}$ 等输入表示的更新，不能归功于 Q/K 学会选择；检验选择能力要用带干扰项的对照实验。
\item[\textbf{C-077.}] \phantomsection\label{solution:C-077}\label{solution:frozen-output-head-training}\hyperref[exercise:frozen-output-head-training]{\textbf{第一步：}}平均交叉熵看三个正确类的概率，贪心预测只看每行最大概率，两者分开记录。\par\nobreak
\textbf{解题过程：}实测损失由 $1.798430$ 降到 $\answerresult{0.736445}$；位置 2 在第 12 步改对，但位置 3 从第 13 步起改猜 ID 5，第 20 步贪心预测为 $[4,5,5]$。平均损失持续下降不保证每个位置都预测正确，需要逐位置检查。
\item[\textbf{C-078.}] \phantomsection\label{solution:C-078}\label{solution:training-rate-transfer}\hyperref[exercise:training-rate-transfer]{\textbf{第一步：}}回想第~\ref*{intro:16a-gradient-descent}~章：步长过大时，每一步越过谷底的幅度会越来越大。\par\nobreak\textbf{解题过程：}$\eta=0.2$ 时第 20 步损失为 $\answerresult{0.035471}$，预测 $[4,5,3]$ 全对；$\eta=1.0$ 时第 5 步已降到 $0.035315$，第 20 步为 $\answerresult{0.003354}$；$\eta=3.0$ 时第 1 步损失升到 $4.080523$，第 2 步 $25.863833$，第 3 步起溢出为 \texttt{inf}，随后变成 \texttt{nan}，$\answerresult{\text{发散}}$。在这次实验、相同初值与给定步数下，$1.0$ 比 $0.2$ 降得更低，$3.0$ 出现损失急升和数值溢出。第~\ref*{intro:16a-gradient-descent}~章算出的 $0<\eta<1$ 只适用于那一个二次损失，不能直接套给本网络；三次运行也不足以确定完整稳定区间。即使两个学习率都能收敛，较大的也可能因振荡而更慢：该章三点例从相同初值训练 10 步，$\eta=0.75$ 的损失约为 $0.500001$，$\eta=0.9$ 则约为 $0.511529$。
\item[\textbf{C-079.}] \phantomsection\label{solution:C-079}\label{solution:training-sentence-change}\hyperref[exercise:training-sentence-change]{\textbf{第一步：}}平均损失是三个位置 $-\ln\Pi_{t,y_t}$ 的均值；先想同一组初始参数给新目标的概率会更高还是更低。\par\nobreak\textbf{解题过程：}初始损失为 $\answerresult{2.436438}$，高于主例：同一组初始参数给新目标的概率更低。20 步后损失降到 $\answerresult{0.746673}$，贪心预测为 $[5,4,4]$：前两个位置已对，第三个位置仍把“花”之后预测成“花”（ID 4），而不是句号。同样的步数和学习率，换一组数据收敛快慢就不同；主例的 20 步结果不能直接搬到别的句子。
\end{enumerate}
\Needspace{6\baselineskip}
\subsubsection*{第~\ref*{intro:23-llm-systems}~章 · 训练好了，模型怎样逐词生成文字？\enspace{\mdseries\footnotesize（\hyperref[exercise:C-080]{题目见第~\pageref*{exercise:C-080}~页起}）}}
\begin{enumerate}[label={},leftmargin=4.8em,itemsep=4pt]
\item[\textbf{C-080.}] \phantomsection\label{solution:C-080}\label{solution:cached-next-token}\hyperref[exercise:cached-next-token]{\textbf{第一步：}}先分开数首次 Prefill 和追加一步的 K/V。\par\nobreak\textbf{解题过程：}在前 $1000$ 个 Token 已完成 Prefill 的前提下，追加第 $1001$ 个后，整段重算需 $1001$ 份 K/V，复用缓存只需新算第 $1001$ 个位置的 $1$ 份 K/V。它节省重复运算，但需要保存各层历史 K/V，内存取决于长度、头数、头维及精度；首次 Prefill 仍要计算。
\item[\textbf{C-081.}] \phantomsection\label{solution:C-081}\label{solution:context-window-memory}\hyperref[exercise:context-window-memory]{\textbf{第一步：}}先检查旧会话文本有没有进入本次上下文。\par\nobreak\textbf{解题过程：}昨天另一会话的内容未加入本次上下文，就不能直接读取。获授权的应用可保存／检索历史摘要并重新放入上下文；推理没有永久改变权重。
\item[\textbf{C-082.}] \phantomsection\label{solution:C-082}\label{solution:greedy-generation}\hyperref[exercise:greedy-generation]{\textbf{第一步：}}每一步只取最后一行概率最大的 ID，并把它接回输入。\par\nobreak
\textbf{解题过程：}实测依次得到 $4$（花，$p=0.8817$）、$5$（开了，$p=0.4566$）、$3$（。，$p=0.5829$），上下文变为 $[2,4,5,3]$。停止是因为教学位置表 $P_{\rm pos}$ 只有 3 行，无法为第 4 个输入位置查表；这是长度上限，不表示模型学会了结束。
\end{enumerate}
\Needspace{6\baselineskip}
\subsubsection*{第~\ref*{intro:23a-from-toy-to-gpt}~章 · 从三维到真实模型：规模、配置与分词\enspace{\mdseries\footnotesize（\hyperref[exercise:C-083]{题目见第~\pageref*{exercise:C-083}~页起}）}}
\begin{enumerate}[label={},leftmargin=4.8em,itemsep=4pt]
\item[\textbf{C-083.}] \phantomsection\label{solution:C-083}\label{solution:uniform-baseline-interpretation}\hyperref[exercise:uniform-baseline-interpretation]{\textbf{第一步：}}先检查初始 logits 是否均匀。\par\nobreak\textbf{解题过程：}不构成代码错误证明；初值给出非均匀 logits，检查标签错位、Softmax、平均方式与参数表，再用可复算脚本核对。
\item[\textbf{C-084.}] \phantomsection\label{solution:C-084}\label{solution:compute-token-budget}\hyperref[exercise:compute-token-budget]{\textbf{第一步：}}先算一层 $12d^2$，再乘层数；宽度出现在平方项里。\par\nobreak
\textbf{解答：}$12\times12\times768^2=84{,}934{,}656$，约 $85$ M。宽度翻倍时每层 $12d^2$ 变为 $\answerresult{4}$ 倍；层数翻倍只使总量变为 $2$ 倍。
\end{enumerate}
\Needspace{6\baselineskip}
\subsubsection*{第~\ref*{intro:23aa-pretraining-icl}~章 · 预训练已经学会续写，为什么还能临时学会新任务？\enspace{\mdseries\footnotesize（\hyperref[exercise:C-085]{题目见第~\pageref*{exercise:C-085}~页起}）}}
\begin{enumerate}[label={},leftmargin=4.8em,itemsep=4pt]
\item[\textbf{C-085.}] \phantomsection\label{solution:C-085}\label{solution:prompt-or-training}\hyperref[exercise:prompt-or-training]{\textbf{第一步：}}先区分格式需求与新季度事实来源。\par\nobreak\textbf{解题过程：}格式变化可先采用 Prompt／Few-shot，大量稳定格式样本可考虑 SFT；读取新季度事实则必须通过授权检索／工具并核对出处。
\item[\textbf{C-086.}] \phantomsection\label{solution:C-086}\label{solution:icl-query-transfer}\hyperref[exercise:icl-query-transfer]{\textbf{第一步：}}先算缩放分数，再做 Softmax 与读出，最后乘输出头。\par\nobreak
\textbf{解题过程：}（a）$s=(0,\ln9)$，$\alpha=(0.1,0.9)$，读出 $(0.1,0.9)$，$Z=(0.5,4.5)$，$\Pr(D)\approx\answerresult{0.982014}$，与查询 $A$ 的情形对称。（b）$s=(\ln3,0)$，$\alpha=(0.75,0.25)$，$Z=(3.75,1.25)$，$\Pr(B)=1/(1+e^{-2.5})\approx\answerresult{0.924142}$。参数没变，查询与示例决定读出；匹配越弱，选择越不尖锐。
\end{enumerate}
\Needspace{6\baselineskip}
\subsubsection*{第~\ref*{intro:23c-post-training}~章 · 怎样改进回答：训练反馈与推理时计算\enspace{\mdseries\footnotesize（\hyperref[exercise:C-087]{题目见第~\pageref*{exercise:C-087}~页起}）}}
\begin{enumerate}[label={},leftmargin=4.8em,itemsep=4pt]
\item[\textbf{C-087.}] \phantomsection\label{solution:C-087}\label{solution:reference-policy-role}\hyperref[exercise:reference-policy-role]{\textbf{第一步：}}看训练是否会更新权重，还是只改变当前上下文。\par\nobreak
\textbf{解答：} SFT 使用示范数据训练；RLVR 利用可验证奖励进行训练；
CoT 提示改变输入或生成形式，本身保持参数固定。
\item[\textbf{C-088.}] \phantomsection\label{solution:C-088}\label{solution:reasoning-interventions}\hyperref[exercise:reasoning-interventions]{\textbf{第一步：}}逐项检查改变的是上下文、参数还是候选选择。\par\nobreak\textbf{解题过程：}Prompt 分步改变本次上下文，SFT 在示例上更新参数，可验证奖励通过可检查结果优化策略；三者都不天然保证每一步推理真实。
\item[\textbf{C-089.}] \phantomsection\label{solution:C-089}\label{solution:candidate-success-budget}\hyperref[exercise:candidate-success-budget]{\textbf{第一步：}}先要求全部失败的概率 $0.7^n\le0.1$。\par\nobreak
\textbf{解题过程：}$n\ge\ln10/\ln(1/0.7)\approx6.456$，最少 $\answerresult{7}$ 份：六份成功率 $1-0.7^6\approx88.2\%$，七份约 $91.8\%$。完全相关时成功率仍只有 $30\%$，复制结果没有独立试错收益。
\item[\textbf{C-090.}] \phantomsection\label{solution:C-090}\label{solution:faulty-verifier-limit}\hyperref[exercise:faulty-verifier-limit]{\textbf{第一步：}}先求“至少一份通过”的概率，再看第一份通过的候选是 42 的条件概率。\par\nobreak\textbf{解题过程：}每份候选不通过的概率是 $1/4$，所以至少一份通过的概率为 $1-0.25^n$；第一份通过的候选在 41、42、43 中等可能，是 42 的概率为 $1/3$。于是 $P_n=(1-0.25^n)/3$：（a）$P_2=0.9375/3=\answerresult{0.3125}$；（b）趋于 $\answerresult{1/3}$；（c）验证器分不出 41、43 与 42，多采样只让“有候选通过”几乎必然发生，却不改变通过者中正确答案只占三分之一。要提高上限，得先改进验证器。
\end{enumerate}
\Needspace{6\baselineskip}
\subsubsection*{第~\ref*{intro:23d-agent-systems}~章 · 检索与证据：RAG 怎样回答有来源的问题？\enspace{\mdseries\footnotesize（\hyperref[exercise:C-091]{题目见第~\pageref*{exercise:C-091}~页起}）}}
\begin{enumerate}[label={},leftmargin=4.8em,itemsep=4pt]
\item[\textbf{C-091.}] \phantomsection\label{solution:C-091}\label{solution:rag-failure-modes}\hyperref[exercise:rag-failure-modes]{\textbf{第一步：}}分别检查检索、来源和生成步骤可能引入的错误。\par\nobreak
\textbf{解答：} 不能。检索可能找错、漏掉或返回过期文档；生成可能误读检索内容或编造未被支持的细节，仍需要校验与来源。
\item[\textbf{C-092.}] \phantomsection\label{solution:C-092}\label{solution:quarterly-evidence-update}\hyperref[exercise:quarterly-evidence-update]{\textbf{第一步：}}先固定相邻季度的分母，再分别算两个增长率。\par\nobreak
\textbf{解题过程：}Q2 环比为 $(125-100)/100=\answerresult{25\%}$；
Q3 环比为 $(132-125)/125=\answerresult{5.6\%}$。
仅修改 Prompt 输入数据不改变网络权重；
实际执行 SFT 训练才会进行参数更新。
\item[\textbf{C-093.}] \phantomsection\label{solution:C-093}\label{solution:rag-query-ranking}\hyperref[exercise:rag-query-ranking]{\textbf{第一步：}}先把查询行向量右乘文档矩阵的转置，再排序。\par\nobreak
\textbf{解题过程：}$q=(0.2,1)$，
$qD^\top=(0.2\cdot1+1\cdot0,\;0.2\cdot0+1\cdot1,\;
0.2\cdot0.8+1\cdot0.8)=(0.2,1,0.96)$。
Top-2 为 Q2、Q3；
Q1 因 $k=2$ 被遗漏。若任务要求三个季度，仍不能据此完成核对。
\item[\textbf{C-094.}] \phantomsection\label{solution:C-094}\label{solution:missing-evidence-response}\hyperref[exercise:missing-evidence-response]{\textbf{第一步：}}先标出缺失证据影响哪一步计算。\par\nobreak\textbf{解题过程：}明确列出“Q2 原始证据缺失”，暂不计算依赖 Q2 的环比。
\end{enumerate}
\Needspace{6\baselineskip}
\subsubsection*{第~\ref*{intro:23d-agent-tools}~章 · 工具调用与 Agent 循环：从请求到执行\enspace{\mdseries\footnotesize（\hyperref[exercise:C-095]{题目见第~\pageref*{exercise:C-095}~页起}）}}
\begin{enumerate}[label={},leftmargin=4.8em,itemsep=4pt]
\item[\textbf{C-095.}] \phantomsection\label{solution:C-095}\label{solution:workflow-agent-choice}\hyperref[exercise:workflow-agent-choice]{\textbf{第一步：}}判断任务是否有稳定步骤，及如何独立复核。\par\nobreak
\textbf{解答：} 示例：固定格式的发票字段抽取可以使用确定工作流并做字段校验；
跨多种格式的资料比较可以用 Agent 选择读取工具，
最终均需核对来源、单位与输出内容。\par\textbf{评分要点：}(1)固定任务写明可预先确定的步骤与格式；(2)动态任务说明为什么要按中间结果选择工具；(3)两者都给出可独立执行的验证，如字段校验、来源核对；(4)指出权限或成本边界。
\item[\textbf{C-096.}] \phantomsection\label{solution:C-096}\label{solution:fixed-workflow-audit}\hyperref[exercise:fixed-workflow-audit]{\textbf{第一步：}}先看步骤和输入格式能否预先确定。\par\nobreak\textbf{解题过程：}固定 Workflow 更易测试、审计与限制成本。\par\textbf{评分要点：}(1)依据是步骤与输入格式能否预先确定；(2)Workflow 可逐步测试、日志可重放；(3)调用次数、成本与权限可预先限定；(4)指出例外：出现未知格式时再引入 Agent 或人工处理。
\end{enumerate}
\Needspace{6\baselineskip}
\subsubsection*{第~\ref*{intro:23d-agent-protocols}~章 · 协议、Skill 与多 Agent：怎样在边界内协作？\enspace{\mdseries\footnotesize（\hyperref[exercise:C-097]{题目见第~\pageref*{exercise:C-097}~页起}）}}
\begin{enumerate}[label={},leftmargin=4.8em,itemsep=4pt]
\item[\textbf{C-097.}] \phantomsection\label{solution:C-097}\label{solution:retrieved-instruction-boundary}\hyperref[exercise:retrieved-instruction-boundary]{\textbf{第一步：}}先区分工具返回数据与用户授权指令。\par\nobreak
\textbf{解答：} 那句话是\textbf{不可信外部内容}而非用户授权指令；
未经适当验证和明确授权，不应执行向陌生外部地址发送文件的动作。
\item[\textbf{C-098.}] \phantomsection\label{solution:C-098}\label{solution:source-permission-boundaries}\hyperref[exercise:source-permission-boundaries]{\textbf{第一步：}}把“排名高”“来源可信”“获得授权”拆成三个独立问题。\par\nobreak
\textbf{解答：}检索排名只反映与查询的相关程度，不等于来源可信或获得授权；网页内容是低信任数据，必须核验来源，并由应用层单独校验工具调用和外发权限。
\item[\textbf{C-099.}] \phantomsection\label{solution:C-099}\label{solution:skill-subagent-design}\hyperref[exercise:skill-subagent-design]{\textbf{第一步：}}先列出 Skill 的交付物与子任务的验收条件。\par\nobreak\textbf{解题过程：}Skill 记录报表字段、单位换算、出处格式和检查脚本；分派给 Subagent 的单位检查任务还需限定可读报告、返回字段及来源，由主任务复核后汇总。若三个文件格式和步骤固定，也可直接运行既定工作流。
\item[\textbf{C-100.}] \phantomsection\label{solution:C-100}\label{solution:a2a-result-verification}\hyperref[exercise:a2a-result-verification]{\textbf{第一步：}}先列出协议传输之外尚未解决的信任问题。\par\nobreak\textbf{解题过程：}A2A 返回成功状态只能表明对方报告了完成；接收方仍核验身份、来源、单位、内容和授权范围。数字无出处时请求原表并标为待核验；回传的越权指令不得充当用户授权。
\end{enumerate}
\Needspace{6\baselineskip}
\subsubsection*{第~\ref*{intro:23b-deployment-and-ecosystem}~章 · 数据、训练系统与部署：资源怎样分配？\enspace{\mdseries\footnotesize（\hyperref[exercise:C-101]{题目见第~\pageref*{exercise:C-101}~页起}）}}
\begin{enumerate}[label={},leftmargin=4.8em,itemsep=4pt]
\item[\textbf{C-101.}] \phantomsection\label{solution:C-101}\label{solution:kv-cache-bytes}\hyperref[exercise:kv-cache-bytes]{\textbf{第一步：}}先按 K/V、层、头、维、位置、字节展开。\par\nobreak\textbf{解题过程：}$2\text{(K/V)}\cdot2\text{层}\cdot2\text{头}\cdot8\text{维}\cdot10\text{位置}\cdot2\text{字节}=\answerresult{1280}$ 字节；KV Cache 随请求长度／并发变化，权重存储由固定模型的参数量与精度决定。
\item[\textbf{C-102.}] \phantomsection\label{solution:C-102}\label{solution:weight-memory-budget}\hyperref[exercise:weight-memory-budget]{\textbf{第一步：}}参数个数乘每个参数的字节数；4 bit 是半个字节。\par\nobreak
\textbf{解答：}fp16 约 $14$ GB；4 bit 约 $\answerresult{3.5}$ GB（未计量化缩放因子等元数据）。部署还需 KV Cache、激活与运行时缓冲，并随上下文长度和并发增长。
\end{enumerate}
\Needspace{6\baselineskip}
\subsubsection*{第~\ref*{intro:22d-modern-transformers}~章 · 现代 Transformer 为什么看起来越来越不像原版？\enspace{\mdseries\footnotesize（\hyperref[exercise:C-103]{题目见第~\pageref*{exercise:C-103}~页起}）}}
\begin{enumerate}[label={},leftmargin=4.8em,itemsep=4pt]
\item[\textbf{C-103.}] \phantomsection\label{solution:C-103}\label{solution:sliding-window-pairs}\hyperref[exercise:sliding-window-pairs]{\textbf{第一步：}}窗口包含当前位置，前两行还未填满窗口。\par\nobreak
\textbf{解题过程：}五行可见位置数为 $(1,2,2,2,2)$，
允许位置对总数 $1+2+2+2+2=\answerresult{9}$；
这里包含查询自己的当前位置。
\item[\textbf{C-104.}] \phantomsection\label{solution:C-104}\label{solution:sliding-window-code}\hyperref[exercise:sliding-window-code]{\textbf{第一步：}}修改窗口宽度，逐行数出允许位置对。\par\nobreak
\textbf{解题过程：}将 $w=3$ 代入五行因果滑窗，
逐行可见位置数为 $(1,2,3,3,3)$，
允许位置对共 $\answerresult{12}$ 个，比 $w=2$ 多 $3$ 个。
可添加 \texttt{assert sum(map(sum, sliding\_mask(5, 3))) == 12}。
\item[\textbf{C-105.}] \phantomsection\label{solution:C-105}\label{solution:gqa-cache-ratio}\hyperref[exercise:gqa-cache-ratio]{\textbf{第一步：}}缓存有 K 和 V 两份，再乘位置数、头数和每头宽度。\par\nobreak
\textbf{解答：} MHA：$2\times120\times6\times4=5{,}760$ 个数；
GQA：$2\times120\times2\times4=1{,}920$ 个数；
每个两字节时是 $3{,}840$ 字节（仅本层这些缓存数值）。
\item[\textbf{C-106.}] \phantomsection\label{solution:C-106}\label{solution:rope-relative-position}\hyperref[exercise:rope-relative-position]{\textbf{第一步：}}Query 旋转 $m\theta$、Key 旋转 $n\theta$；点积只依赖两者之间的夹角。\par\nobreak
\textbf{解答：} 相对角为 $(5-2)\theta=3\theta$。
\item[\textbf{C-107.}] \phantomsection\label{solution:C-107}\label{solution:moe-routing}\hyperref[exercise:moe-routing]{\textbf{第一步：}}对分数排序即可确定 Top-1 与 Top-2。\par\nobreak
\textbf{解答：} Top-1 选专家 2；Top-2 选专家 2 和 3。
每次只激活路由选中的专家。
\item[\textbf{C-108.}] \phantomsection\label{solution:C-108}\label{solution:architecture-fair-comparison}\hyperref[exercise:architecture-fair-comparison]{\textbf{第一步：}}先固定质量、数据、硬件和负载条件。\par\nobreak\textbf{解题过程：}固定数据、质量指标、硬件、精度、负载和输入／生成长度，联合比较准确性、显存、延迟、吞吐及成本。\par\textbf{评分要点：}(1)每次只改一个模块，其余配置相同；(2)质量与资源指标同时报告；(3)负载、长度、精度与硬件固定并写明；(4)多次运行给出波动或置信区间。
\end{enumerate}
\Needspace{6\baselineskip}
\subsubsection*{第~\ref*{intro:23e-tech-judgement}~章 · 面对新模型与新缩写，怎么判断它值不值得学？\enspace{\mdseries\footnotesize（\hyperref[exercise:C-109]{题目见第~\pageref*{exercise:C-109}~页起}）}}
\begin{enumerate}[label={},leftmargin=4.8em,itemsep=4pt]
\item[\textbf{C-109.}] \phantomsection\label{solution:C-109}\label{solution:production-architecture-evaluation}\hyperref[exercise:production-architecture-evaluation]{\textbf{第一步：}}先写出双方共用的测试条件表。\par\nobreak\textbf{解题过程：}固定任务、精度、硬件、输入／输出长度及并发；共同报告质量、延迟、吞吐、显存与成本，分析失败案例与统计波动。\par\textbf{评分要点：}(1)至少四项可复现的控制条件；(2)同时报告质量与延迟、吞吐、显存、成本；(3)说明榜单可能受污染或与业务分布不同；(4)给出失败案例或统计波动。
\item[\textbf{C-110.}] \phantomsection\label{solution:C-110}\label{solution:benchmark-count-gap}\hyperref[exercise:benchmark-count-gap]{\textbf{第一步：}}用 $\sqrt{p(1-p)/n}$ 估计标准误，再与差值比较。\par\nobreak
\textbf{解题过程：}$\sqrt{0.81\times0.19/500}\approx0.0175$，约 $1.8$ 个百分点；若按两份独立估计近似，差值的标准误约 $2.5$ 个百分点，$0.8$ 远小于它，不能断言前者更好。还应做同题配对比较、报告多次运行的波动，并排查评测污染与提示格式差异。
\end{enumerate}
\Needspace{6\baselineskip}
\subsubsection*{第~\ref*{intro:19-cnn}~章 · CNN：怎样用局部连接与参数共享处理图片？\enspace{\mdseries\footnotesize（\hyperref[exercise:C-111]{题目见第~\pageref*{exercise:C-111}~页起}）}}
\begin{enumerate}[label={},leftmargin=4.8em,itemsep=4pt]
\item[\textbf{C-111.}] \phantomsection\label{solution:C-111}\label{solution:convolution-output-shape}\hyperref[exercise:convolution-output-shape]{\textbf{第一步：}}用 $(4-2)/1+1$ 求每个空间方向的输出尺寸。\par\nobreak
\textbf{解答：} 宽高为 $3\times3$；卷积核权重 $2\times2=4$ 个（不含偏置）。
\item[\textbf{C-112.}] \phantomsection\label{solution:C-112}\label{solution:convolution-shared-kernel}\hyperref[exercise:convolution-shared-kernel]{\textbf{第一步：}}分别取左上与右上两个 $2\times2$ 窗口，各与同一个核对应相乘。\par\nobreak
\textbf{解题过程：}左上输出为 $1\times1+2\times0+0\times0+1\times(-1)=\answerresult{0}$；右上输出为 $2\times1+0\times0+1\times0+3\times(-1)=\answerresult{-1}$。两个位置共用一组 $4$ 个核权重，不是各有一组；输入窗口不同，输出可以不同。
\item[\textbf{C-113.}] \phantomsection\label{solution:C-113}\label{solution:pooling-detail-loss}\hyperref[exercise:pooling-detail-loss]{\textbf{第一步：}}先看下采样保留了哪些空间信息。\par\nobreak\textbf{解题过程：}池化可降低空间尺寸与后续成本，但可能丢失精细位置、边缘或小目标信息。
\end{enumerate}
\Needspace{6\baselineskip}
\subsubsection*{第~\ref*{intro:19b-vit}~章 · 从 CNN 到 ViT：如何把图片变成一串 Token？\enspace{\mdseries\footnotesize（\hyperref[exercise:C-114]{题目见第~\pageref*{exercise:C-114}~页起}）}}
\begin{enumerate}[label={},leftmargin=4.8em,itemsep=4pt]
\item[\textbf{C-114.}] \phantomsection\label{solution:C-114}\label{solution:vit-spatial-order}\hyperref[exercise:vit-spatial-order]{\textbf{第一步：}}先问平均后是否还能定位原来的区域。\par\nobreak\textbf{解题过程：}无序平均抹掉空间身份，难定位左上区域；可加入二维位置编码或保留 Patch 序列及相应位置标识。
\item[\textbf{C-115.}] \phantomsection\label{solution:C-115}\label{solution:patch-projection-count}\hyperref[exercise:patch-projection-count]{\textbf{第一步：}}先数行、列方向各能切几块，再把一块的高、宽、通道数相乘。\par\nobreak
\textbf{解答：}$(32/8)^2=16$ 个 Patch；每块展平为 $8\times8\times3=192$；投影矩阵 $192\times64$，共 $\answerresult{12{,}288}$ 个权重，所有 Patch 共用；加 CLS 后序列长度为 $17$。
\end{enumerate}
\Needspace{6\baselineskip}
\subsubsection*{第~\ref*{intro:04-key-value-outer-product}~章 · 线性联想记忆：写入、串扰与误差修正\enspace{\mdseries\footnotesize（\hyperref[exercise:C-116]{题目见第~\pageref*{exercise:C-116}~页起}）}}
\begin{enumerate}[label={},leftmargin=4.8em,itemsep=4pt]
\item[\textbf{C-116.}] \phantomsection\label{solution:C-116}\label{solution:memory-single-write}\hyperref[exercise:memory-single-write]{\textbf{第一步：}}先用 Key 的每个分量分别乘整条 Value，形成状态矩阵的两行。\par\nobreak
\textbf{解题过程：}$\mathcal{M}=\begin{bmatrix}1.2&-0.6\\1.6&-0.8\end{bmatrix}$；读取结果为 $(0.6\times1.2+0.8\times1.6,\,-0.6\times0.6-0.8\times0.8)=\answerresult{(2,-1)}$。一般有 $\vect k^\top\mathcal{M}=(\vect k^\top\vect k)\vect v^\top$；单位 Key 的自身点积为 $1$，所以读回 $\vect v^\top$。
\item[\textbf{C-117.}] \phantomsection\label{solution:C-117}\label{solution:memory-cross-talk}\hyperref[exercise:memory-cross-talk]{\textbf{第一步：}}将查询对两条外积的读取分别展开，再相加。\par\nobreak
\textbf{解题过程：}$\mathcal{M}=\begin{bmatrix}1&0.6\\0&0.8\end{bmatrix}$，$\vect k_A^\top\mathcal{M}=\answerresult{(1,0.6)}$，多出的 $(0,0.6)$ 来自 B。一般交叉项是 $(\vect k_A^\top\vect k_B)\vect v_B^\top$；非零交叉点积把别条 Value 混进读出。
\item[\textbf{C-118.}] \phantomsection\label{solution:C-118}\label{solution:memory-outer-shape}\hyperref[exercise:memory-outer-shape]{\textbf{第一步：}}先写出 $\vect{v}^\top$ 与 $\vect{k}$ 的形状，比较内侧维度。\par\nobreak
\textbf{解题过程：}$\vect{v}^\top:1\times3$ 乘 $\vect{k}:2\times1$，内侧维度 $3\neq2$，无法相乘；正确写入 $\vect{k} \vect{v}^\top$ 为 $2\times3$，供 $\vect{q}^\top:1\times2$ 读出 $1\times3$。若 $d_k=d_v$，$\vect{v}^\top\vect{k}$ 是 $1\times1$ 标量（两向量的内积），不是 $d_k\times d_v$ 的状态矩阵。
\item[\textbf{C-119.}] \phantomsection\label{solution:C-119}\label{solution:delta-write-order}\hyperref[exercise:delta-write-order]{\textbf{第一步：}}依次写入 B、计算 A 的残差、再重读旧 B。\par\nobreak
\textbf{解答：} 先写 B 得
$\mathcal{M}_1=\vect{k}_B \vect{v}_B^\top=\begin{bmatrix}0&0.6\\0&0.8\end{bmatrix}$；
再写 A 时，$\vect{k}_A^\top\mathcal{M}_1=(0,0.6)$，残差
$\vect{r}_A^\top=(1,-0.6)$，因此
$\mathcal{M}_2=\mathcal{M}_1+\vect{k}_A \vect{r}_A^\top=
\boxed{\begin{bmatrix}1&0\\0&0.8\end{bmatrix}}$。
最新的 A 读取为 $\vect{k}_A^\top\mathcal{M}_2=\answerresult{(1,0)}$；
旧的 B 读取为 $\vect{k}_B^\top\mathcal{M}_2=\answerresult{(0.6,0.64)}$，
误差 $\sqrt{0.6^2+(-0.36)^2}\approx\answerresult{0.700}$。
与正文先 A 后 B 对称：\textbf{最后写入的 Key 当次精确，旧 Key 被扰动}，
不能推断两条同时精确。
\item[\textbf{C-120.}] \phantomsection\label{solution:C-120}\label{solution:memory-write-strategies}\hyperref[exercise:memory-write-strategies]{\textbf{第一步：}}先检查修正 B 时会不会同时改变 A 的读出。\par\nobreak
\textbf{解题过程：}接近正交会减小不同 Key 的交叉点积；Delta 写入则按当前读出与目标 Value 的差补写。非正交时，修正新记忆也可能扰动旧记忆，因此单遍修正不保证同时读准。可以监测全部读出误差并尝试反复校正；若 Key 与目标允许同一个线性状态精确表示，也可求联立方程。反复校正的收敛仍取决于更新规则与步长。
\end{enumerate}
\Needspace{6\baselineskip}
\subsubsection*{第~\ref*{intro:09-nonlinearity-and-memory}~章 · 非线性检索：从 MLP 视角到记忆容量\enspace{\mdseries\footnotesize（\hyperref[exercise:C-121]{题目见第~\pageref*{exercise:C-121}~页起}）}}
\begin{enumerate}[label={},leftmargin=4.8em,itemsep=4pt]
\item[\textbf{C-121.}] \phantomsection\label{solution:C-121}\label{solution:memory-pairing-shape}\hyperref[exercise:memory-pairing-shape]{\textbf{第一步：}}先核对权重列数与 Value 行数。\par\nobreak\textbf{解题过程：}四条 Key 会得到 $1\times4$ 匹配权重，而三条 Value 只有三行，$1\times4$ 与 $3\times4$ 不能相乘。
\item[\textbf{C-122.}] \phantomsection\label{solution:C-122}\label{solution:memory-selectivity}\hyperref[exercise:memory-selectivity]{\textbf{第一步：}}先比较幂运算后的弱／强比值；Softmax 则比较两个指数的商。\par\nobreak
\textbf{解题过程：}幂运算后的弱／强比值依次为 $\answerresult{0.5,0.25,0.125}$，强匹配仍为 $1$。Softmax 的第一条与第二条权重之比为 $e^{\beta(1-0.6)}$，从 $e^{0.4}$ 增至 $e^{0.8}$，因此第一条权重上升。这两种方法都让已有分数的差异更突出；本题的幂运算比较限于 $[0,1]$ 内的非负分数。
\item[\textbf{C-123.}] \phantomsection\label{solution:C-123}\label{solution:memory-relu-read}\hyperref[exercise:memory-relu-read]{\textbf{第一步：}}先让 $\vect{x}^\top$ 分别与 $K$ 的两行做点积，再逐项过 ReLU。\par\nobreak
\textbf{解答：}$\vect{x}^\top K^\top=(1-2,\,3-1)=(-1,2)$，所以 $\relu(\vect{x}^\top K^\top)=\answerresult{(0,2)}$：负匹配的第 1 条记忆被关闭，第 2 条按分数 $2$ 保留。
\item[\textbf{C-124.}] \phantomsection\label{solution:C-124}\label{solution:memory-mlp-interpretation}\hyperref[exercise:memory-mlp-interpretation]{\textbf{第一步：}}先区分对计算机制的解释与对输出事实的验证。\par\nobreak
\textbf{解题过程：}这种对应把各隐藏分量的激活和输出贡献联系起来，帮助理解网络怎样组合特征。它并未建立单个神经元与一条真实事实的一一对应。确认回答仍要核查原始来源、语境与实际输出；单个激活不是事实证据。
\item[\textbf{C-125.}] \phantomsection\label{solution:C-125}\label{solution:memory-random-capacity}\hyperref[exercise:memory-random-capacity]{\textbf{第一步：}}先用正交向量条数不超过维度，区分精确上限与随机量级估计。\par\nobreak
\textbf{解题过程：}最多 $4$ 条；前四条若两两正交，第 $5$ 条非零 Key 必与其中某条有非零点积。随机交叉点积的典型量级分别为 $\answerresult{0.125}$ 与 $\answerresult{0.03125}$。应用于真实模型前，应检查真实 Key 的相关性是否符合随机假设，并测随记忆数变化的读出误差及 Value 尺度；维度量级不能直接解释成能记住的可靠事实数。
\end{enumerate}
\Needspace{6\baselineskip}
\subsubsection*{附录 A · 优化器\enspace{\mdseries\footnotesize（\hyperref[exercise:C-126]{题目见第~\pageref*{exercise:C-126}~页起}）}}
\begin{enumerate}[label={},leftmargin=4.8em,itemsep=4pt]
\item[\textbf{C-126.}] \phantomsection\label{solution:C-126}\label{solution:optimizer-clipping}\hyperref[exercise:optimizer-clipping]{\textbf{第一步：}}先求长度 $\sqrt{5^2+12^2}$，只有超过阈值才整体缩放。\par\nobreak
\textbf{解题过程：}$\|(5,12)\|_2=13$，乘 $6.5/13=0.5$ 得 $\answerresult{(2.5,6)}$，方向不变；阈值为 $20$ 时长度未超限，梯度保持 $(5,12)$。梯度裁剪作用于更新前的梯度，不是对预测概率归一化。高级优化器手算见附录 A。
\item[\textbf{C-127.}] \phantomsection\label{solution:C-127}\label{solution:optimizer-warmup}\hyperref[exercise:optimizer-warmup]{\textbf{第一步：}}每一步先用上一步的 $w$ 重新算梯度，再乘该步的学习率。\par\nobreak
\textbf{解题过程：}$w_1=1-0.01\times2=\answerresult{0.98}$，$w_2=0.98(1-0.04)=\answerresult{0.9408}$，$w_3=0.9408(1-0.06)\approx\answerresult{0.8844}$。恒定 $\eta=0.03$ 时 $w_3=0.94^3\approx0.8306$，下降更快。Warmup 只改变步长 $\eta_t$，不更换损失函数；初期步子小，前几步下降也较慢。
\item[\textbf{C-128.}] \phantomsection\label{solution:C-128}\label{solution:optimizer-adamw}\hyperref[exercise:optimizer-adamw]{\textbf{第一步：}}第一步偏置校正后 $\widehat{m}_1/\sqrt{\widehat{v}_1}=g_1/|g_1|$，先算出这一方向。\par\nobreak
\textbf{解题过程：}（a）方向为 $1$，$w_1=(1-0.1\times0.05)\times4-0.1\times1=\answerresult{3.88}$。（b）$g_1'=0.5+0.2=0.7$，但 $\widehat{m}_1/\sqrt{\widehat{v}_1}=0.7/0.7=1$，$w_1=4-0.1=\answerresult{3.9}$：衰减在第一步被自适应归一化抵消。AdamW 把缩小项放在梯度统计之外，才保证每步都按 $\eta\lambda w$ 衰减。
\end{enumerate}
\Needspace{6\baselineskip}
\subsubsection*{累积复习 · 学习\enspace{\mdseries\footnotesize（\hyperref[exercise:R-01]{题目见第~\pageref*{exercise:R-01}~页起}）}}
\begin{enumerate}[label={},leftmargin=4.8em,itemsep=4pt]
\item[\textbf{R-01.}] \phantomsection\label{solution:R-01}\label{solution:review-learning-cycle}\hyperref[exercise:review-learning-cycle]{\textbf{第一步：}}先按时间顺序列出哪些数被计算、哪些数被改变。\par\nobreak\textbf{解题过程：}前向产生预测，损失比较目标，反传给出旧参数处的梯度，优化器统一更新。验证集检验新样本表现，不能把训练损失下降等同于泛化。
\item[\textbf{R-02.}] \phantomsection\label{solution:R-02}\label{solution:review-validation-experiment}\hyperref[exercise:review-validation-experiment]{\textbf{第一步：}}先冻结测试集，再排查划分泄漏和模型选择。\par\nobreak\textbf{解题过程：}固定数据划分、随机种子与预算，对比一次正则化或容量调整；用验证集选模，最终测试集只用于独立报告。记录两条损失曲线与所改因素。
\end{enumerate}
\Needspace{6\baselineskip}
\subsubsection*{累积复习 · 矩阵\enspace{\mdseries\footnotesize（\hyperref[exercise:R-03]{题目见第~\pageref*{exercise:R-03}~页起}）}}
\begin{enumerate}[label={},leftmargin=4.8em,itemsep=4pt]
\item[\textbf{R-03.}] \phantomsection\label{solution:R-03}\label{solution:review-matrix-shapes}\hyperref[exercise:review-matrix-shapes]{\textbf{第一步：}}先固定样本轴，再按权重改变特征轴。\par\nobreak\textbf{解题过程：}$X:3\times2$，$W_1:2\times4$，激活 $3\times4$，$W_2:4\times2$，输出 $3\times2$；每行用同一函数，输入不同仍可得到不同输出。
\item[\textbf{R-04.}] \phantomsection\label{solution:R-04}\label{solution:review-matrix-nonlinearity}\hyperref[exercise:review-matrix-nonlinearity]{\textbf{第一步：}}先用矩阵乘法结合律，再选含负中间值的输入。\par\nobreak\textbf{解题过程：}无非线性时两式相等（浮点计算允许舍入差）；可取标量 $X=-1,W_1=W_2=1$，线性输出 $-1$，ReLU 后输出 $0$。
\end{enumerate}
\Needspace{6\baselineskip}
\subsubsection*{累积复习 · 文字\enspace{\mdseries\footnotesize（\hyperref[exercise:R-05]{题目见第~\pageref*{exercise:R-05}~页起}）}}
\begin{enumerate}[label={},leftmargin=4.8em,itemsep=4pt]
\item[\textbf{R-05.}] \phantomsection\label{solution:R-05}\label{solution:review-token-pipeline}\hyperref[exercise:review-token-pipeline]{\textbf{第一步：}}先错开一位构造目标，再区分词表轴与特征轴。\par\nobreak\textbf{解题过程：}完整 ID $[2,4,5,3]$；输入 $[2,4,5]$、目标 $[4,5,3]$。三维 Embedding 后 $3\times3$，概率 $3\times5$，第一位置目标用 \texttt{probs[0,3]}。
\item[\textbf{R-06.}] \phantomsection\label{solution:R-06}\label{solution:review-cross-entropy}\hyperref[exercise:review-cross-entropy]{\textbf{第一步：}}使用 $\ln(ab)=\ln a+\ln b$。\par\nobreak\textbf{解题过程：}惊讶度减少 $\ln2$。两位置若为 $1/4,1/2$，联合概率 $1/8$，负对数 $\ln8=\ln4+\ln2$。
\textbf{一种可行方案：}例如用标准库 \texttt{math.log} 分别计算 \texttt{-math.log(p*q)} 与 \texttt{-math.log(p)-math.log(q)}，断言差的绝对值小于 1e-12；概率必须为正。
\end{enumerate}
\Needspace{6\baselineskip}
\subsubsection*{累积复习 · 上下文\enspace{\mdseries\footnotesize（\hyperref[exercise:R-07]{题目见第~\pageref*{exercise:R-07}~页起}）}}
\begin{enumerate}[label={},leftmargin=4.8em,itemsep=4pt]
\item[\textbf{R-07.}] \phantomsection\label{solution:R-07}\label{solution:review-context-evidence}\hyperref[exercise:review-context-evidence]{\textbf{第一步：}}检查两条样本的当前位置是否相同。\par\nobreak\textbf{解题过程：}两例当前词与位置均相同，位置标签也无法区分。RNN 可把前词编码进状态，Attention 可直接读取前词；关键是输出收到不同前文。
\item[\textbf{R-08.}] \phantomsection\label{solution:R-08}\label{solution:review-multihead-shapes}\hyperref[exercise:review-multihead-shapes]{\textbf{第一步：}}匹配轴是位置，拼接轴是特征。\par\nobreak\textbf{解题过程：}每头分数 $3\times3$，可见位置数为 $1,2,3$；Value 为 $3\times3$，读出 $3\times3$，两头拼接 $3\times6$，输出投影再决定残差宽度。
\textbf{一种可行方案：}示例在窗口外放“春天／冬天”，窗口内均为相同填充词与“花”；目标分别“开／谢”。固定窗口看到相同输入，能读取早先线索的方法可区分；RNN 也可传递线索。
\end{enumerate}
\Needspace{6\baselineskip}
\subsubsection*{累积复习 · Transformer\enspace{\mdseries\footnotesize（\hyperref[exercise:R-09]{题目见第~\pageref*{exercise:R-09}~页起}）}}
\begin{enumerate}[label={},leftmargin=4.8em,itemsep=4pt]
\item[\textbf{R-09.}] \phantomsection\label{solution:R-09}\label{solution:review-training-generation}\hyperref[exercise:review-training-generation]{\textbf{第一步：}}先找训练标签，再画生成反馈箭头。\par\nobreak\textbf{解题过程：}训练有错位目标，因果 Mask 防止偷看，反传更新参数。生成只有前缀，选出的 Token 成为下一输入，参数固定；遇结束 ID 或长度等条件停止。
\item[\textbf{R-10.}] \phantomsection\label{solution:R-10}\label{solution:review-cache-output-head}\hyperref[exercise:review-cache-output-head]{\textbf{第一步：}}检查缓存对应的输入、位置与参数是否匹配。\par\nobreak\textbf{解题过程：}Prefill 处理已知前缀并写 K/V；Decode 只为新 Token 写入并读取历史。改变前文会改变其后各层状态，不能无条件复用原缓存；只可按实现复用完全匹配的前缀。
\textbf{一种可行方案：}按 \path{teaching_reference.py} 的 \texttt{parameters/forward/backward} 顺序计算，保存 forward 返回的 loss 与 A 的第三行，统一更新；冻结 Wout 时只跳过它的更新，其他梯度仍反传，按相同初值和步数比较。
\end{enumerate}
\Needspace{6\baselineskip}
\subsubsection*{累积复习 · 助手\enspace{\mdseries\footnotesize（\hyperref[exercise:R-11]{题目见第~\pageref*{exercise:R-11}~页起}）}}
\begin{enumerate}[label={},leftmargin=4.8em,itemsep=4pt]
\item[\textbf{R-11.}] \phantomsection\label{solution:R-11}\label{solution:review-assistant-boundaries}\hyperref[exercise:review-assistant-boundaries]{\textbf{第一步：}}把证据缺失与指令权限分开处理。\par\nobreak\textbf{解题过程：}补检并核验 Q2，按同一单位和口径重算环比；网页指令不构成用户授权，拒绝越权外发。检索、工具和 Prompt 改动本身不更新权重，训练才更新。
\item[\textbf{R-12.}] \phantomsection\label{solution:R-12}\label{solution:review-deployment-cost}\hyperref[exercise:review-deployment-cost]{\textbf{第一步：}}先分开计算阶段与内存项。\par\nobreak\textbf{解题过程：}固定硬件、精度、输入/输出长度及并发，报告首 Token 延迟、后续速度、峰值显存、质量与总成本；模型权重之外还有缓存、运行时缓冲，多采样增加生成和验证开销。
\textbf{一种可行方案：}至少列模型版本、量化、硬件、输入/输出长度、批量与并发；同时测首 Token 延迟、每 Token 延迟、吞吐、峰值显存与质量，证据核验另计成本。
\end{enumerate}
\Needspace{6\baselineskip}
\subsubsection*{累积复习 · 架构\enspace{\mdseries\footnotesize（\hyperref[exercise:R-13]{题目见第~\pageref*{exercise:R-13}~页起}）}}
\begin{enumerate}[label={},leftmargin=4.8em,itemsep=4pt]
\item[\textbf{R-13.}] \phantomsection\label{solution:R-13}\label{solution:review-gqa-budget}\hyperref[exercise:review-gqa-budget]{\textbf{第一步：}}先数单层、单请求的 K/V 元素，再乘层数、请求数和每元素字节数。\par\nobreak
\textbf{解题过程：}GQA 单层、单请求有 $2\times100\times3\times8=4{,}800$ 个数；两层、两条请求的缓存本体为 $4{,}800\times2\times2\times2=\answerresult{38{,}400}$ 字节。改为 MHA 后，K/V 头数从 $3$ 变为 $12$，字节数为 $38{,}400\times(12/3)=\answerresult{153{,}600}$。倍率是 $\answerresult{4}$，请求数同时乘在两者上，不改变倍率；未计布局、分配与运行时开销。
\par\textbf{真实配置：}层数为 $24$，隐藏宽度为 $896$，Q/KV 头数分别为 $14/2$。头宽由隐藏宽度除以 Q 头数得到 $896/14=\answerresult{64}$；缓存计数使用 KV 头数 $2$。单位置、全层缓存为 $2\times24\times2\times64\times2=12{,}288$ 字节；两条请求各 $100$ 个位置，共 $12{,}288\times100\times2=\answerresult{2{,}457{,}600}$ 字节。MHA 有 $14$ 个 KV 头，倍率为 $14/2=\answerresult{7}$。\texttt{scripts/config\_table.py} 可离线核对四个字段；\texttt{scripts/cache\_budget\_demo.py} 使用同一固定配置，复算每位置 $12{,}288$ 字节并推广到 32K。
\item[\textbf{R-14.}] \phantomsection\label{solution:R-14}\label{solution:review-architecture-ablation}\hyperref[exercise:review-architecture-ablation]{\textbf{第一步：}}每次只替换一个组件并保持负载一致。\par\nobreak\textbf{解题过程：}固定任务、权重训练预算、硬件、精度、长度、并发与质量阈值，分别替换位置、KV 共享、内核，再测组合。联合报告质量、延迟、吞吐与显存，不能用一个总时间归因三项。
\end{enumerate}
\Needspace{6\baselineskip}
\subsubsection*{累积复习 · 视觉\enspace{\mdseries\footnotesize（\hyperref[exercise:R-15]{题目见第~\pageref*{exercise:R-15}~页起}）}}
\begin{enumerate}[label={},leftmargin=4.8em,itemsep=4pt]
\item[\textbf{R-15.}] \phantomsection\label{solution:R-15}\label{solution:review-image-token-shapes}\hyperref[exercise:review-image-token-shapes]{\textbf{第一步：}}Patch 个数决定行数，投影决定列数。\par\nobreak\textbf{解题过程：}$4\times4$ 的 Patch 行矩阵，乘 $4\times3$ 投影得 $4\times3$；乘 $3\times4$ 桥接得 $4\times4$，与 $2\times4$ 文字拼成 $6\times4$。
\item[\textbf{R-16.}] \phantomsection\label{solution:R-16}\label{solution:review-vision-sharing}\hyperref[exercise:review-vision-sharing]{\textbf{第一步：}}先区分局部特征提取与跨位置读取。\par\nobreak\textbf{解题过程：}卷积核在不同窗口复用；Patch 投影对各 Patch 用同一矩阵。ViT 的 Attention 跨 Patch 读取；卷积通过局部窗口与多层感受野传播。
\textbf{一种可行方案：}可用同一四乘四图片仅改变左上二乘二块。Patch 线性投影首先改变对应行；卷积首先改变感受野覆盖该像素的输出；全局 Attention 后影响可传播到其他位置。
\end{enumerate}
\Needspace{6\baselineskip}
\subsubsection*{累积复习 · 联想记忆\enspace{\mdseries\footnotesize（\hyperref[exercise:R-17]{题目见第~\pageref*{exercise:R-17}~页起}）}}
\begin{enumerate}[label={},leftmargin=4.8em,itemsep=4pt]
\item[\textbf{R-17.}] \phantomsection\label{solution:R-17}\label{solution:review-memory-pipeline}\hyperref[exercise:review-memory-pipeline]{\textbf{第一步：}}用分配律展开，再检查归一化分母。\par\nobreak\textbf{解题过程：}$\vect{q}^\top\mathcal{M}=\sum_i(\vect{q}^\top\vect{k}_i)\vect{v}_i^\top$；二者都按匹配量加权 Value，但 Softmax 对当前 Query 逐行归一化，外积线性状态没有该分母，不能直接替换。
\item[\textbf{R-18.}] \phantomsection\label{solution:R-18}\label{solution:review-memory-correction}\hyperref[exercise:review-memory-correction]{\textbf{第一步：}}把新写入外积左乘旧 Key。\par\nobreak\textbf{解题过程：}变化量含 $(\vect{k}_A^\top\vect{k}_B)\vect{r}_B^\top$，交叉点积非零会改变旧读出。逐步增加实际 Key/Value 数，记录相关性、读取误差与失败比例，固定维度及读取规则。
\end{enumerate}
